\chapter{Introducción}

\section{Propósito}

Este manual explica cómo instalar DIGAIPE en un servidor de la red institucional, de principio
a fin, de dos maneras:

\begin{itemize}
	\item \textbf{Con contenedores} (capítulo~\ref{cap:docker}): todos los componentes se
	      ejecutan en contenedores definidos en \cod{deploy/compose.yaml}, con Docker o con Podman
	      ejecutado por un usuario sin privilegios (sección~\ref{sec:podman}). En este manual,
	      ``Docker'' designa este modo con cualquiera de los dos.
	\item \textbf{En servidor físico} (sin Docker), en Linux (capítulo~\ref{cap:linux}) o en
	      Windows (capítulo~\ref{cap:windows}): cada componente se instala como un servicio del
	      sistema operativo.
\end{itemize}

Ambas instalaciones dan el mismo resultado: la misma base de datos, las mismas cuentas con
permisos mínimos, HTTPS con una autoridad certificadora local, respaldos semanales y un registro
de auditoría verificable. El servidor descarga lo que necesita (código, paquetes e imágenes)
directamente de internet durante la instalación y las actualizaciones.

\section{Componentes}

\begin{xltabular}{\linewidth}{@{}l L L@{}}
	\toprule
	\textbf{Componente} & \textbf{Función} & \textbf{Implementación} \\
	\midrule
	\endhead
	Base de datos & Datos de admisiones, usuarios, sesiones y registro de auditoría &
	MySQL 8.4 LTS \\
	Aplicación & Interfaz web y lógica de negocio & Node.js 24 (SvelteKit, \texttt{node build}),
	escucha solo en el propio servidor, puerto 3000 \\
	Proxy HTTPS & Única entrada desde la red. Emite y renueva el certificado solo &
	Caddy 2 con autoridad certificadora local (\texttt{local\_certs}) o, si el servidor es
	público, Let's Encrypt \\
	Instalación (\emph{setup}) & Migraciones, triggers de auditoría, cuentas de MySQL,
	administrador inicial y datos base. Se ejecuta y termina &
	\texttt{scripts/docker/setup.sh} y los scripts de \texttt{scripts/} \\
	Respaldos & Copia semanal de la base y de los archivos cargados, con suma SHA-256 &
	\texttt{deploy/backup/*.sh} (bash, \texttt{mysqldump}) \\
	Auditoría & Verificación y anclaje de la cadena de hashes &
	\texttt{scripts/audit-ledger.ts} \\
	\bottomrule
\end{xltabular}

\begin{codigo}[label={Flujo de red (igual en ambos modos)}]
  PCs de la red
       |  HTTPS 443  (80 solo redirige a HTTPS)
       v
  +---------+   HTTP 127.0.0.1:3000   +------------+   TCP 3306   +-----------+
  |  Caddy  | ----------------------> | Aplicación | -----------> | MySQL 8.4 |
  +---------+                         +------------+              +-----------+
                                            |                          ^
                                     var/uploads (archivos)            |
                                            |                          |
                                            v                          |
                                     +-----------+  mysqldump          |
                                     | Respaldos | --------------------+
                                     +-----------+
                                            |
                                            v
                                  carpeta de respaldos (otro disco)
\end{codigo}

Ni MySQL (3306) ni la aplicación (3000) aceptan conexiones desde la red: solo Caddy publica
puertos. Esto no es opcional; la aplicación confía en la cabecera \cod{X-Forwarded-For} que
añade Caddy para registrar la dirección IP real de cada usuario en la auditoría. Si la
aplicación fuera accesible directamente, cualquiera podría falsificar esa dirección.

\section{¿Docker o servidor físico?}

\begin{xltabular}{\linewidth}{@{}l L L@{}}
	\toprule
	& \textbf{Docker} & \textbf{Servidor físico} \\
	\midrule
	\endhead
	Ventaja principal & Es la instalación que mantiene el repositorio
	(\cod{deploy/}); un solo archivo de configuración; instalar y actualizar son dos o tres
	comandos &
	No requiere Docker; funciona en Windows Server; cada pieza se ve y se administra con
	las herramientas habituales del sistema \\
	Desventaja & Requiere Docker. En Windows, Docker Desktop tiene limitaciones serias para
	un servidor (sección~\ref{sec:docker-windows}) & Más pasos, más archivos de configuración;
	las actualizaciones se hacen a mano \\
	Respaldos & Servicio \cod{backup-worker} incluido & Servicio \texttt{systemd} (Linux) o
	tarea programada (Windows) \\
	Dificultad & Baja & Media \\
	\bottomrule
\end{xltabular}

\begin{nota}[Recomendación]
	\textbf{Use Docker sobre Linux} siempre que pueda: es la configuración probada por el equipo de
	desarrollo y la que menos pasos manuales tiene. Use la instalación en servidor físico cuando la
	política institucional no permita Docker o cuando el servidor disponible sea Windows Server.
\end{nota}

\section{Convenciones de este manual}

\begin{itemize}
	\item \so{Linux} Los comandos se ejecutan en una terminal con un usuario que puede usar
	      \cod{sudo}.
	\item \so{Windows} Los comandos se ejecutan en \textbf{PowerShell como administrador}
	      (clic derecho → Ejecutar como administrador), salvo que se indique otra cosa.
	\item Lo que va entre \texttt{<} y \texttt{>} se reemplaza por su valor, sin los signos.
	      Por ejemplo, \texttt{<DIGAIPE\_HOST>} se reemplaza por \texttt{digaipe.ula.local}.
	\item Las rutas son literales y las mismas en todo el manual. Si elige otras, cámbielas en
	      \emph{todos} los archivos y comandos. El apéndice~\ref{ap:ubicaciones} las resume.
\end{itemize}

\begin{aviso}[Datos personales]
	La base de datos, la carpeta de archivos cargados y los respaldos contienen datos de
	estudiantes, hashes de contraseñas y sesiones. Los archivos de configuración contienen las
	contraseñas de MySQL. Trátelos todos como información confidencial: no los envíe por correo,
	no los suba a servicios externos y no los deje en carpetas compartidas sin control de acceso.
\end{aviso}
